\vspace{4pt plus 4pt minus 1pt}\noindent{\kbvsecfont\fontsize{11pt}{14.5pt}\selectfont\color{kbvaccentdark} (೩)\hspace{0.45em}ಅಲ್ಪ ವಿರಾಮ ಚಿಹ್ನೆ:}\hspace{0.5em}({\kbvlatin i}) ಸಂಬೋಧನಾರ್ಥಕ ಶಬ್ದದ ಮುಂದೆ, ({\kbvlatin ii}) ಕರ್ತೃ, ಕರ್ಮ, ಕ್ರಿಯಾಪದಗಳಿಗೆ ಒಂದಕ್ಕಿಂತ ಹೆಚ್ಚು ವಿಶೇಷಣಗಳಿರುವಾಗ, ({\kbvlatin iii}) ಒಂದೇ ಕ್ರಿಯೆಗೆ ಒಂದಕ್ಕಿಂತ ಹೆಚ್ಚು ಕರ್ತೃ ಪದಗಳಿರುವಾಗ, ({\kbvlatin iv}) ಒಂದೇ ಕ್ರಿಯೆಗೆ ಒಂದಕ್ಕಿಂತ ಹೆಚ್ಚು ಕರ್ಮಪದ ಗಳಿರುವಾಗ ಅಲ್ಪ ವಿರಾಮ ಚಿಹ್ನೆಯನ್ನು ಕೆಳಗೆ ಕೊಟ್ಟಿರುವ ಉದಾಹರಣೆ ಗಳಲ್ಲಿ ತೋರಿಸಿರುವಂತೆ ಬಳಸಬೇಕು.\par

\vspace{3pt plus 2pt}\noindent\hspace*{1.2em}\kbvlabel{ಉದಾ}\par\vspace{1pt}

\vspace{2.5pt plus 3pt minus 0.5pt}\noindent\hspace*{1.2em}\kbvnum{(1)}\hspace{0.55em}ಪದ್ಮಲತಾ, ಬೀಗದ ಕೈ ಜೋಪಾನವಾಗಿಡು.\par

\vspace{2.5pt plus 3pt minus 0.5pt}\noindent\hspace*{1.2em}\kbvnum{(2)}\hspace{0.55em}ಕನ್ಯೆಯು ರೂಪವತಿಯೂ, ಗುಣವತಿಯೂ, ವಿದ್ಯಾವತಿಯೂ ಆಗಿದ್ದಾಳೆ.\par

\vspace{2.5pt plus 3pt minus 0.5pt}\noindent\hspace*{1.2em}\kbvnum{(3)}\hspace{0.55em}ನಮ್ಮ ಅಂಗಡಿಯಲ್ಲಿ ಶುದ್ಧವಾದ, ಹಸನಾದ, ಮಂಡ್ಯದ ಬೆಣ್ಣೆ ಮಾರಾಟಕ್ಕಿದೆ.\par

\vspace{2.5pt plus 3pt minus 0.5pt}\noindent\hspace*{1.2em}\kbvnum{(4)}\hspace{0.55em}ಚಪಾತಿಯು ಬಿಸಿಯಾಗಿಯೂ, ಮೃದುವಾಗಿಯೂ, ರುಚಿ ಯಾಗಿಯೂ ಇರಬೇಕು.\par

\vspace{2.5pt plus 3pt minus 0.5pt}\noindent\hspace*{1.2em}\kbvnum{(5)}\hspace{0.55em}ಧರ್ಮರಾಯ, ಭೀಮ, ಅರ್ಜುನ, ನಕುಲ, ಸಹದೇವ ಇವರೈ ವರೂ ದ್ರೌಪದಿಯನ್ನು ವಿವಾಹವಾದರು.\par

\vspace{2.5pt plus 3pt minus 0.5pt}\noindent\hspace*{1.2em}\kbvnum{(6)}\hspace{0.55em}ನಮ್ಮ ಕೃಷಿ ಕ್ಷೇತ್ರದಲ್ಲಿ ಕಬ್ಬು, ತೆಂಗು, ಭತ್ತ, ರಾಗಿ, ತೊಗರಿ ಬೆಳೆಯುತ್ತೇವೆ.\par

\vspace{4pt plus 4pt minus 1pt}\noindent{\kbvsecfont\fontsize{11pt}{14.5pt}\selectfont\color{kbvaccentdark} (೪)\hspace{0.45em}ಪ್ರಶ್ನಾರ್ಥಕ ಚಿಹ್ನೆ:}\hspace{0.5em}ಪ್ರಶ್ನಾರ್ಥಕ ಶಬ್ದದ ಮುಂದೆ, ಪ್ರಶ್ನಾರ್ಥಕ ವಾಕ್ಯದ ಮುಂದೆ ಈ ಚಿಹ್ನೆಯನ್ನು ಬಳಸಬೇಕು.\par

\noindent\hspace*{1.2em}\kbvlabel{ಉದಾ}\hspace{0.6em}ರೈಲು ಈ ದಿನ ತಡವಾಗಿ ಬಂದಿತು ; ಏಕೆ?\par

\noindent\hspace*{2.2em}ನೀನು ಈ ದಿನ ಎಲ್ಲಿಗೆ ಹೋಗುತ್ತೀಯೆ ?\par

\vspace{4pt plus 4pt minus 1pt}\noindent{\kbvsecfont\fontsize{11pt}{14.5pt}\selectfont\color{kbvaccentdark} (೫)\hspace{0.45em}ಭಾವಸೂಚಕ ಚಿಹ್ನೆ:}\hspace{0.5em}ಭಾವಸೂಚಕ ಶಬ್ದದ ಮುಂದೆ, ಭಾವಸೂಚಕ ವಾಕ್ಯದ ಮುಂದೆ ಈ ಚಿಹ್ನೆಯನ್ನು ಬಳಸಬೇಕು.\par
